% Options for packages loaded elsewhere
\PassOptionsToPackage{unicode}{hyperref}
\PassOptionsToPackage{hyphens}{url}
\PassOptionsToPackage{dvipsnames,svgnames,x11names}{xcolor}
%
\documentclass[
  10pt,
  a4paper,
]{article}
\usepackage{amsmath,amssymb}
\usepackage{setspace}
\usepackage{iftex}
\ifPDFTeX
  \usepackage[T1]{fontenc}
  \usepackage[utf8]{inputenc}
  \usepackage{textcomp} % provide euro and other symbols
\else % if luatex or xetex
  \usepackage{unicode-math} % this also loads fontspec
  \defaultfontfeatures{Scale=MatchLowercase}
  \defaultfontfeatures[\rmfamily]{Ligatures=TeX,Scale=1}
\fi
\usepackage{lmodern}
\ifPDFTeX\else
  % xetex/luatex font selection
  \setmainfont[]{Liberation Serif}
  \setsansfont[]{Liberation Sans}
  \setmonofont[]{DejaVu Sans Mono}
\fi
% Use upquote if available, for straight quotes in verbatim environments
\IfFileExists{upquote.sty}{\usepackage{upquote}}{}
\IfFileExists{microtype.sty}{% use microtype if available
  \usepackage[]{microtype}
  \UseMicrotypeSet[protrusion]{basicmath} % disable protrusion for tt fonts
}{}
\makeatletter
\@ifundefined{KOMAClassName}{% if non-KOMA class
  \IfFileExists{parskip.sty}{%
    \usepackage{parskip}
  }{% else
    \setlength{\parindent}{0pt}
    \setlength{\parskip}{6pt plus 2pt minus 1pt}}
}{% if KOMA class
  \KOMAoptions{parskip=half}}
\makeatother
\usepackage{xcolor}
\usepackage[margin=22mm]{geometry}
\setlength{\emergencystretch}{3em} % prevent overfull lines
\providecommand{\tightlist}{%
  \setlength{\itemsep}{0pt}\setlength{\parskip}{0pt}}
\setcounter{secnumdepth}{-\maxdimen} % remove section numbering
\ifLuaTeX
\usepackage[bidi=basic]{babel}
\else
\usepackage[bidi=default]{babel}
\fi
\babelprovide[main,import]{korean}
\ifPDFTeX
\else
\babelfont{rm}[]{Liberation Serif}
\fi
% get rid of language-specific shorthands (see #6817):
\let\LanguageShortHands\languageshorthands
\def\languageshorthands#1{}
\usepackage{xeCJK}
\setCJKmainfont{Noto Serif CJK KR}
\setCJKsansfont{Noto Sans CJK KR}
\setCJKmonofont{Noto Sans Mono CJK KR}
\xeCJKsetup{CJKspace=true}
\usepackage{xurl}
\usepackage{fvextra}
\usepackage{fancyhdr}
\AtBeginDocument{\DefineVerbatimEnvironment{verbatim}{Verbatim}{breaklines=true,breakanywhere=true,fontsize=\small}}
\setlength{\emergencystretch}{4em}
\setlength{\headheight}{15pt}
\pagestyle{fancy}\fancyhf{}
\fancyhead[L]{\small 중성자별 구조와 비방사 진동}
\fancyfoot[C]{\thepage}
\renewcommand{\headrulewidth}{0.3pt}
\setcounter{tocdepth}{2}
\ifLuaTeX
  \usepackage{selnolig}  % disable illegal ligatures
\fi
\usepackage{bookmark}
\IfFileExists{xurl.sty}{\usepackage{xurl}}{} % add URL line breaks if available
\urlstyle{same}
\hypersetup{
  pdftitle={A20 · EOS의 값과 미분을 같은 구간에서 읽기},
  pdflang={ko-KR},
  colorlinks=true,
  linkcolor={Maroon},
  filecolor={Maroon},
  citecolor={Blue},
  urlcolor={Blue},
  pdfcreator={LaTeX via pandoc}}

\title{A20 · EOS의 값과 미분을 같은 구간에서 읽기}
\usepackage{etoolbox}
\makeatletter
\providecommand{\subtitle}[1]{% add subtitle to \maketitle
  \apptocmd{\@title}{\par {\large #1 \par}}{}{}
}
\makeatother
\subtitle{과제 큐시트}
\author{}
\date{}

\begin{document}
\maketitle

\setstretch{1.12}
\textbf{선행:} A13, A18　 \textbf{독립 작업:} 2--4시간 (계획치)

\subsection{같은 표에서 읽은 값과 미분이 왜 서로
어긋날까?}\label{uxac19uxc740-uxd45cuxc5d0uxc11c-uxc77duxc740-uxac12uxacfc-uxbbf8uxbd84uxc774-uxc65c-uxc11cuxb85c-uxc5b4uxae0buxb0a0uxae4c}

압력과 첫 미분이 연속인 cubic을 이어 붙였다. 둘째 미분도 연속인가?
압축응답의 knot를 물질적 밀도 jump로 읽어도 되는지 먼저 적자.

먼저 한두 문장으로 답한 뒤 아래 설명을 읽자. 첫 예측이 틀렸거나 아직
모르겠다고 적는 것은 감점하지 않는다.

\subsection{필요한 관계와 읽을
위치}\label{uxd544uxc694uxd55c-uxad00uxacc4uxc640-uxc77duxc744-uxc704uxce58}

\(n=P_{,t}/H,\ \epsilon=\mu n-P,\ \chi=H P_{,t}/(\mu P_{,tt}),\ dr/dx=H/(\mu\nu')\).

\begin{itemize}
\tightlist
\item
  \path|notes/chapters/ch09.md|
\item
  \path|methods/src/eos/ns_cell_eos.f90|
\item
  \path|methods/src/eos/ns_table_eos.f90|
\item
  \path|calculation/tests_fortran/cell_checks.f90|
\end{itemize}

\subsection{계산으로
옮기기}\label{uxacc4uxc0b0uxc73cuxb85c-uxc62euxae30uxae30}

제공 함수의 호출이 기본 경로다. 구간 보간·역함수·적응 적분·matching의
내부 코드를 작성하지 않고, 입력 조건과 반환된 근거를 해석한다.

\begin{enumerate}
\def\labelenumi{\arabic{enumi}.}
\tightlist
\item
  P와 첫 미분의 연속성을 둘째 미분의 연속성과 구분한다.
\item
  무차원 P=mu\^{}2와 좁은 구간 예제에서 제공 cell\_state·역함수를
  호출한다.
\item
  calculation의 cell\_checks를 실행하고 12개 관계가 의미하는 범위를
  적는다.
\item
  다른 t와 구간 폭에서 제1법칙과 SI 차원을 손계산으로 확인한다.
\end{enumerate}

\newpage

\subsection{실행해서 무엇을 볼
것인가}\label{uxc2e4uxd589uxd574uxc11c-uxbb34uxc5c7uxc744-uxbcfc-uxac83uxc778uxac00}

패키지 루트에서 실행한다. \texttt{-\/-plan}은 경로만 확인하며 계산
성공을 뜻하지 않는다.

\begin{verbatim}
python tools/run_assignment.py A20 --plan
python tools/run_assignment.py A20 --run --output ../A20.json
\end{verbatim}

\subsection{결과를
설명하기}\label{uxacb0uxacfcuxb97c-uxc124uxba85uxd558uxae30}

평가할 대상은 하나의 구간 함수다. Native knot를 보존하고 작은 local
coordinate에서 값·미분·역함수를 확인한다. 표를 더 조밀하게 만든 것과는
다르다.

\textbf{남길 증거.} 한 구간의 P,n,epsilon,chi 유도; local t 역함수와
절대 mu subtraction 비교; 선택 workspace의 실제 검사 기록

\textbf{판단 기준.} 2점: 계수 jump와 새로운 물질적 jump를 구분하고, 각
미분의 H 차수와 pressure residual을 설명한다. 최초 예측의 정오보다,
결과를 본 뒤 이유를 고치고 다른 입력에도 적용하는지를 본다.

\subsection{막혔을 때}\label{uxb9c9uxd614uxc744-uxb54c}

\textbf{힌트 1.} 미분할 때 t의 단위와 H 인자를 확인한다.

\textbf{힌트 2.} P-a0와 t(a1+t(a2+a3t))를 비교하면 작은 차이를 절대 mu로
되돌리지 않아도 된다.

\subsection{다음 만남에 다시
묻기}\label{uxb2e4uxc74c-uxb9ccuxb0a8uxc5d0-uxb2e4uxc2dc-uxbb3buxae30}

다음 회차에는 구간 폭만 다른 같은 시험문제를 고른다. 재보간 없이 역함수
잔차를 비교하고, 정밀도가 나빠졌다면 어떤 연산을 먼저 의심할지 말한다.
먼저 짧게 답한 뒤 본문과 자신의 기록을 열어 확인한다.

\textbf{범위.} Cell 검사는 별의 고유값 또는 실제 EOS의 오차 상한을
증명하지 않는다.

\end{document}
